% Page 046 — page imprimée 42
% Contenu seul : le préambule, l'en-tête et la pagination sont gérés par meyrueis.tex

menaçant, de la figure du garnement, au risque de provoquer un accident. Les enfants
se tenaient tranquilles ; terrifiés par le concierge. Plusieurs, devenus vieux
aujourd’hui, en gardent le souvenir ! Ils me l’ont confirmé il y a peu de temps !
\emph{(je revois encore, en 1940, le Rascle nous poursuivant dans le temple avec son tisonnier, car nous le
faisions enrager ; chauffé à blanc ? bien sûr que non ! se souvient Paul Loupiac)}

Quant aux baptêmes, mariages et enterrements, il exigeait que les familles, après
s’être entretenues avec le pasteur, aillent le trouver, et surtout ne l’oublient pas sous
peine de ne pouvoir accéder au temple, qu’il estimait être un peu sa propriété
personnelle. En cas de négligence, il menaçait de prendre le large à l’heure de la
cérémonie, et de se rendre avec les clefs dans sa poche dans son champ le plus
éloigné. Cela arriva au moins une fois et n’amusa personne !!!

Irène tenait habituellement l’harmonium. Mais, avant et après la naissance d’André,
elle dut renoncer provisoirement à cette charge. Elle fut remplacée par une jeune fille
pleine de bonne volonté, mais qui ne pouvait accompagner les cantiques que d’une
main. L’instrument était bon, un harmonium de quatre jeux et demi, tout à fait adapté
à l’édifice. Un dimanche où je mettais au point avec un groupe de fidèles et
l’organiste, la partie musicale, Mr. Avesque s’approcha de nous et dit :
« \emph{Mademoiselle Agulhon, vous êtes beaucoup plus forte que le pasteur. Lui il a
besoin de deux mains pour jouer, vous, vous vous en tirez avec un doigt !} » On
pouvait considérer qu’il s’agissait d’une réflexion amusante dont on aurait dû
seulement sourire, mais cela déplut visiblement à l’organiste !

A cette période vivait à Meyrueis un pasteur à la retraite, originaire du pays, tout à
fait digne et remarquable. Il se nommait David Poujol.* Le concierge avait usé avec
lui ses fonds de culotte sur les bancs de l'école communale de Meyrueis. Ce pasteur
avait dirigé des églises importantes, Mazamet, Nîmes, La Haye et d’autres lieux.
Orateur incontestable, doué d’une puissante voix, on le sentait habitué aux nombreux
rassemblements. Il cultivait l’éloquence classique à laquelle il excellait. Un
dimanche, alors qu’il montait en chaire (où il me remplaçait, accompagné du
concierge qui le tutoyait depuis toujours, ce dernier lui dit devant moi : « \emph{Alors
David, c’est toi qui vas faire l’article aujourd’hui ?} » ce qui lui attira une verte
réplique, méritée, qui lui cloua le bec !
Je désire rendre hommage à ce pasteur distingué, disparu depuis longtemps, et qui
imposait le respect. Il finissait ses jours à Meyrueis où il vivait dans sa modeste
maison veuf et seul. Au gros de l’hiver il se rendait chez une de ses filles mais
revenait très vite. C’était un ramasseur de champignons. [\ldots]

Puisque je parle de David Poujol, comment ne pas mentionner le rayonnement de
deux autres pasteurs qui passaient, à cette époque leurs vacances d’été à Meyrueis.
Emile Guiraud de Paris et Louis Hubac de Castres. Ce dernier était le fils d’un paysan
que nous appréciions beaucoup, habitant la vallée de la Brèze, propriétaire de la
ferme de Conillergues près de Campis. Lorsqu’il prêchait, le temple se remplissait.
Sa prédication originale ( je me souviens d’un sermon sur « la main ») attirait, mais
n’était-il pas, surtout, un enfant du pays ?

\begin{footnotesize}
* David Poujol était l’oncle de Marcel Poujol dit « Poujol d’Alger », rue saint blaise, décédé en 1994.
\end{footnotesize}
